\documentclass[11pt]{article}

\usepackage[margin=1in]{geometry}
\usepackage[T1]{fontenc}
\usepackage[utf8]{inputenc}
\usepackage{lmodern}
\usepackage{microtype}
\usepackage{amsmath}
\usepackage{amssymb}
\usepackage{array}
\usepackage{booktabs}
\usepackage{graphicx}
\usepackage{caption}
\usepackage{placeins}
\usepackage[hidelinks]{hyperref}
\usepackage[authoryear]{natbib}

\setlength{\parindent}{0pt}
\setlength{\parskip}{0.6em}

\newcommand{\code}[1]{\nolinkurl{#1}}
\newcommand{\repositoryurl}{\url{https://github.com/yuhaoyang266/vi-insertion-only-sim}}
\newcommand{\mm}{\ensuremath{\mathrm{mm}}}
\newcommand{\mps}{\ensuremath{\mathrm{m/s}}}
\newcommand{\N}{\ensuremath{\mathrm{N}}}
\graphicspath{{../figures/main/}{../figures/appendix/}{../supplement/figures/}}

\newcommand{\includeorplaceholder}[2]{%
  \IfFileExists{../figures/main/#1.pdf}{%
    \includegraphics[width=\linewidth]{../figures/main/#1.pdf}%
  }{%
    \IfFileExists{../figures/main/#1.png}{%
      \includegraphics[width=\linewidth]{../figures/main/#1.png}%
    }{%
      \IfFileExists{../figures/appendix/#1.pdf}{%
        \includegraphics[width=\linewidth]{../figures/appendix/#1.pdf}%
      }{%
        \IfFileExists{../figures/appendix/#1.png}{%
          \includegraphics[width=\linewidth]{../figures/appendix/#1.png}%
        }{%
          \IfFileExists{../supplement/figures/#1.pdf}{%
            \includegraphics[width=\linewidth]{../supplement/figures/#1.pdf}%
          }{%
            \IfFileExists{../supplement/figures/#1.png}{%
              \includegraphics[width=\linewidth]{../supplement/figures/#1.png}%
            }{%
              \fbox{\parbox[c][0.22\textheight][c]{0.95\linewidth}{%
                \centering
                \textbf{#2}\\[0.5em]
                \small Expected asset under \texttt{figures/} or \texttt{supplement/figures/}: \texttt{\detokenize{#1.pdf}} or \texttt{\detokenize{#1.png}}%
              }}%
            }%
          }%
        }%
      }%
    }%
  }%
}

\title{Support-Gated Variable-Impedance Learning in a 3DoF Insertion Benchmark}
\author{Yu Haoyang\\Shanghai University\\Shanghai Institute of Applied Mathematics and Mechanics\\\texttt{yuhaoyang@shu.edu.cn}}
\date{\small Project repository: \repositoryurl}

\begin{document}

\maketitle

\begin{abstract}
We study Support-Gated Variable-Impedance Learning (SG-VI) in a simulation-only 3DoF final-stage
insertion benchmark with a structured variable-impedance action space
$(\Delta x, \Delta y, \Delta z, K_{xy}, K_z)$. In this paper SG-VI denotes a benchmark-local
recipe that combines explicit variable-impedance commands, behavior-cloning warm-start, and
factorized support controls over demonstrations, resets, and PPO budget. To make the gate
measurable, we define a Support Coverage Index (SCI), a quantized rollout-to-demonstration
state-action overlap diagnostic under the same relative-frame interface. Under a matched
cross-family pure-RL pilot, PPO, SAC, and TD3 from scratch remain outside useful contact from 50k
to 200k environment steps, so optimizer family alone does not reopen the gate. A factorized
diagnostic sweep identifies BC demonstration support as the cleanest lever: removing it drives
useful-contact success to zero, whereas reset coverage, BC depth, and PPO budget show
protocol-sensitive and partly non-monotone effects. A separate pure-clearance Sprint 4A ladder
isolates \code{clearance\_easy}, \code{nominal}, and \code{clearance\_hard} without changing
offset, friction, or force-noise axes; on this ladder \code{BC-only} stays at 1.000 success,
\code{BC$\rightarrow$PPO} and the learned fixed-impedance baseline remain near 0.8 success without
jam, while \code{DAPG-lite} drops to 0.768 success with 0.05 jam under
\code{clearance\_hard}. Variable impedance retains localized physical value under high-friction
contact: tuned fixed stiffness can nearly saturate success, but same-interface mechanics traces
show lower contact load and work for variable stiffness. The claim is benchmark-local and
teacher-coupled, not a general robotics or teacher-independent
algorithm-ranking result.
\end{abstract}

\section{Introduction}

Contact-rich insertion is often framed as a reinforcement-learning problem, but in practice
learning can fail before policy optimization has any chance to matter. A policy may never enter
recoverable contact, or its demonstrations may fail to cover the approach states from which
contact-relevant behavior becomes reachable. In that setting, saying that ``PPO fails'' is too
coarse. The more precise question is whether the system ever provides a learnable interface in the
first place.

This paper studies that question in a simulation-only 3DoF final-stage peg-in-hole benchmark with
a structured variable-impedance action space. The benchmark is intentionally reduced: it models
translational insertion only, exposes the contact transition directly, and uses a policy interface
that separates translational intent from impedance modulation. This reduced scope is a strength for
our purpose. It lets us isolate where the real bottleneck lies: in demonstration guidance, in
upstream reset coverage, in contact-law design, or in the downstream choice among PPO-, SAC-,
TD3-, and DAPG-style optimization.

The study is therefore framed around evidence layers rather than a single algorithm leaderboard.
The benchmark separates whether a policy ever reaches recoverable contact, whether demonstrations
and reset curricula expose useful state-action support, whether PPO-style fine-tuning changes a
BC-initialized policy, and whether variable stiffness changes the contact load path once contact is
reached. We also separate the frozen five-profile benchmark from later stress ladders such as the
pure-clearance Sprint 4A sweep, because those ladders answer narrower robustness questions rather
than replacing the main benchmark contract. This leads to a narrow claim: in the matched
teacher-coupled setup, learnability is dominated by upstream interface design, while the physical
value of variable impedance is localized to lower-load high-friction contact rather than broad
algorithmic superiority.

We package that interface-level view as Support-Gated Variable-Impedance Learning (SG-VI): a
benchmark-local recipe that keeps impedance explicit in the action space, uses BC to reopen
contact-relevant support, and then studies what policy optimization changes once that support is
available. To keep the gate measurable rather than only rhetorical, we also define a Support
Coverage Index (SCI), a discretized rollout-to-demo support-overlap diagnostic under the same
relative-frame interface.

Our contributions are therefore fivefold. First, we propose Support-Gated Variable-Impedance
Learning (SG-VI), a benchmark-local methodology for auditing 3DoF insertion learnability through
explicit impedance commands, BC warm-start, and factorized support controls. Second, we propose
Support Coverage Index (SCI), a benchmark-local discretized state-action support diagnostic that
measures what fraction of a policy rollout remains inside the demonstration support set. Third, we
show with a matched cross-family pure-RL pilot that PPO, SAC, and TD3 without demonstrations do
not reach useful contact under the nominal-train 50k--200k contract, while the mixed-contract
evidence matrix keeps this negative result separate from the five-profile demo-supported benchmark.
Fourth, we show that BC-based support is the decisive gate under the matched teacher-coupled
setup, while a frozen teacher-support-quality $\times$ demo-budget boundary keeps the next teacher
ablation auditable rather than post hoc. Fifth, we show that robustness beyond the frozen
five-profile benchmark remains localized: the main learned fixed-impedance baseline underperforms
primarily in the \code{high\_friction} evaluation class, while a separate pure-clearance ladder
preserves \code{BC-only} success but exposes the clearest harder-clearance penalty in
\code{DAPG-lite}.

The remainder of the paper is organized as follows. Section~\ref{sec:related-work} positions the
study relative to contact-rich insertion, impedance control, and demonstration-guided RL.
Section~\ref{sec:setup} defines the task, action space, comparison set, and evaluation protocol.
Section~\ref{sec:experiments} presents the main benchmark, the cross-family pure-RL negative
evidence, the factorized and teacher-boundary support analysis, the clearance-shift sweep, and the
physical interpretation of the fixed-versus-variable comparison. Section~\ref{sec:discussion}
discusses the scope of the claim and its limitations. The appendix expands the supporting evidence
with the PPO-only large-budget audit, the tuned fixed-stiffness check, evaluation-class
transparency, raw per-profile tables, the contact-transition audit, and the BC stability sweep
behind the reported setting.

\section{Related Work}
\label{sec:related-work}

Peg-in-hole assembly has long served as a canonical contact-rich manipulation problem because small
alignment errors can rapidly produce large contact forces, wedging, and jamming. Recent reviews
continue to organize the field around compliant search, force-feedback correction, and contact
handling rather than around purely kinematic tracking \citep{li2025_peg_review}. Our paper shares
that focus on contact-rich behavior, but it asks a narrower upstream question: before discussing
policy quality, when is the insertion interface learnable at all?

Impedance control provides the classical language for that question. Hogan's original formulation
made explicit that manipulation under interaction cannot be reduced to pure position or pure force
control \citep{hogan1985_impedance_part1,hogan1985_impedance_part2}. Hybrid position/force and
compliance-control formalisms likewise emphasize that contact tasks require explicit task-frame
reasoning and force-motion decomposition rather than only Cartesian tracking
\citep{mason1981_compliance_force,raibert1981_hybrid_position_force}. More recent variable-impedance
work has shown that impedance can itself be learned or adapted from demonstrations, inverse
reinforcement learning, or direct policy optimization
\citep{abudakka2020_variable_impedance_demo,zhang2021_variable_impedance_irl,
zhang2023_variable_impedance_passivity,zhou2024_mobile_manipulator_imitation,
yang2025_variable_impedance_passivity}. We follow that structured-control tradition, but we do not
claim a new variable-impedance controller. Instead, we use an explicit impedance action space to
make learnability bottlenecks auditable.

Demonstration-guided learning is the second anchor. Learning from demonstration has long been used
to provide task-relevant priors when exploration is brittle or unsafe
\citep{argall2009_lfd_survey,billard2008_robot_programming_by_demonstration,ijspeert2013_dmp}.
Demonstration-augmented policy gradients and related imitation-plus-RL pipelines have been
particularly influential in dexterous and contact-rich manipulation
\citep{vecerik2017_demo_sparse_rewards,rajeswaran2018_dexterous_demo_rl,
johannink2019_residual_rl,wang2022_adaptive_imitation_insertion}. Off-policy and offline RL
baselines such as SAC, benchmark suites such as robosuite, and offline datasets and algorithms
such as D4RL and CQL define a broader comparison space that this paper does not exhaust
\citep{haarnoja2018_sac,zhu2020_robosuite,fu2020_d4rl,kumar2020_cql}. Our emphasis is
different from the usual sample-efficiency story. In the present benchmark, demonstrations are not
merely a data-efficiency boost; they are a precondition for entering the contact-relevant state
distribution at all.

\section{Task Setup and Experimental Protocol}
\label{sec:setup}

\subsection{Learnability Question and Evidence Boundary}

We study final-stage translational insertion in an analytical 3DoF environment. The state contains
relative peg position, translational velocity, a synthetic three-axis force signal, and the
previous action. The policy acts in a structured 5D action space
$a_t = (\Delta x, \Delta y, \Delta z, K_{xy}, K_z)$, where the first three dimensions specify local
Cartesian motion and the last two specify lateral and axial stiffness commands. This design keeps
motion and compliance explicit rather than hiding them inside torque-level outputs.

Figure~\ref{fig:overview} summarizes the task geometry and the control interface. The peg is
observed in a hole-centered relative frame, so the policy receives relative pose, relative
velocity, and force information rather than absolute world coordinates. This relative-frame design
is important later: it strongly attenuates the effect of \code{offset\_bias} under the current
environment semantics, although we do not yet isolate that interpretation with a separate
absolute-versus-relative observation ablation.

\begin{figure}[t]
  \centering
  \includeorplaceholder{fig1_task_policy_impedance_overview}{Final Figure 1: task, policy, and impedance overview}
  \caption{3DoF insertion task, policy interface, and impedance strategy. The policy observes the
  peg in a hole-centered relative frame and outputs translational motion commands together with
  lateral and axial stiffness commands. The distinction between fixed and variable impedance is not
  an abstract controller label: it is whether stiffness remains constant during contact or is
  reduced after contact onset.}
  \label{fig:overview}
\end{figure}

\subsection{Mechanics and Interface Formalization}

The policy operates in a hole-centered relative frame. At time step $t$, the observation is
\begin{equation}
\mathbf{o}_t =
\left[\mathbf{p}^{\mathrm{rel}}_t,\; \dot{\mathbf{p}}^{\mathrm{rel}}_t,\; \mathbf{f}_t,\;
\mathbf{a}_{t-1}\right],
\end{equation}
where $\mathbf{p}^{\mathrm{rel}}_t \in \mathbb{R}^3$ is the peg position relative to the hole
frame, $\dot{\mathbf{p}}^{\mathrm{rel}}_t \in \mathbb{R}^3$ is the translational velocity,
$\mathbf{f}_t \in \mathbb{R}^3$ is the synthetic force observation returned by the analytical
contact model, and $\mathbf{a}_{t-1} \in \mathbb{R}^5$ is the previous action. The resulting
observation dimension is 14. The action is
\begin{equation}
\mathbf{a}_t =
\left[\Delta x_t,\; \Delta y_t,\; \Delta z_t,\; \kappa^{xy}_t,\; \kappa^{z}_t\right],
\end{equation}
where the first three dimensions specify local Cartesian motion increments and the last two are
normalized stiffness commands.

The environment decodes the stiffness channels linearly after clipping them to $[0,1]$:
\begin{align}
K_{xy,t} &= K_{xy}^{\min} +
\mathrm{clip}(\kappa^{xy}_t, 0, 1)\left(K_{xy}^{\max} - K_{xy}^{\min}\right), \\
K_{z,t} &= K_{z}^{\min} +
\mathrm{clip}(\kappa^{z}_t, 0, 1)\left(K_{z}^{\max} - K_{z}^{\min}\right),
\end{align}
with the structured stiffness matrix
\begin{equation}
\mathbf{K}_t = \mathrm{diag}(K_{xy,t}, K_{xy,t}, K_{z,t}).
\end{equation}

The analytical contact model returns a synthetic force observation
$\mathbf{f}_t = \mathcal{C}(\mathbf{p}^{\mathrm{rel}}_t, \dot{\mathbf{p}}^{\mathrm{rel}}_t,
\mathbf{K}_t; \theta_{\mathrm{env}})$. When the peg is already inside the hole but still rubbing
the wall, the model applies an in-hole drag term whose lateral component is proportional to
$K_{xy,t}$ and a wall-proximity factor $\rho_t$, while its axial component scales with
$\mu_{\mathrm{wall}} K_{z,t} |\min(\Delta z_t, 0)|$. When lateral overflow creates blocked contact,
the blocked-contact branch uses
\begin{align}
\mathbf{f}^{xy}_{t,\mathrm{block}} &=
\alpha_{xy} K_{xy,t} o_t \mathbf{u}_t + 0.25\, \mu_{\mathrm{wall}} f^0_{z,t} \mathbf{u}_t, \\
f^0_{z,t} &= \alpha_z K_{z,t} |\min(\Delta z_t, 0)|, \\
f_{z,t,\mathrm{block}} &= f^0_{z,t} + 0.2 \|\mathbf{f}^{xy}_{t,\mathrm{block}}\|_2,
\end{align}
where $o_t$ is the lateral overflow magnitude and $\mathbf{u}_t$ is the corresponding unit
direction. Here $\alpha_{xy}$ and $\alpha_z$ are fixed lateral and axial contact-scaling
constants of the analytical environment, $\mu_{\mathrm{wall}}$ is the active wall-friction
coefficient, and $\rho_t$ is the wall-proximity factor returned by the analytical contact
geometry. The final force observation is the sum of the active drag and blocked-contact terms and
is also used for jam monitoring.

A rollout is counted as successful when
\begin{equation}
\|\mathbf{p}^{\mathrm{rel}}_{t,xy}\|_2 \le 0.8\,\mm, \qquad
|p^{\mathrm{rel}}_{t,z}| \le 1.0\,\mm, \qquad
\|\dot{\mathbf{p}}^{\mathrm{rel}}_t\|_2 \le 0.08~\mathrm{m/s},
\end{equation}
and it is terminated as jam when
\begin{equation}
\|\mathbf{f}_t\|_2 \ge 8.0~\N
\end{equation}
for three consecutive steps. The episode horizon is 64 steps. Table~\ref{tab:env_params}
summarizes the key benchmark constants.

\begin{table}[t]
  \centering
  \small
  \caption{Key environment and termination parameters used throughout the main benchmark. The
  stiffness ranges are the default environment decoding ranges; the learned fixed-impedance
  baseline replaces them with constant midpoint values during both training and evaluation.}
  \label{tab:env_params}
  \begin{tabular}{ll}
    \toprule
    Parameter & Value \\
    \midrule
    Observation dimension & 14 \\
    Action dimension & 5 \\
    $K_{xy}$ range & $[20.0,\;110.0]$ \\
    $K_z$ range & $[35.0,\;140.0]$ \\
    Success tolerance & lateral $\le 0.8\,\mm$, axial $\le 1.0\,\mm$ \\
    Success speed gate & $\|\dot{\mathbf{p}}^{\mathrm{rel}}\|_2 \le 0.08~\mathrm{m/s}$ \\
    Jam criterion & $\|\mathbf{f}\|_2 \ge 8.0~\N$ for 3 consecutive steps \\
    Episode horizon & 64 steps \\
    \bottomrule
  \end{tabular}
\end{table}

\subsection{Compared Pipelines}

The reported comparison set contains five methods.
\begin{enumerate}
  \item \textbf{PPO w/o BC}: direct policy optimization from scratch.
  \item \textbf{BC-only}: a behavior-cloned policy evaluated without PPO fine-tuning.
  \item \textbf{Fixed-impedance RL}: the same BC-based learning path, but with stiffness channels
  pinned to constant values.
  \item \textbf{BC$\rightarrow$PPO}: the mainline that warm-starts PPO from the stable BC base.
  \item \textbf{DAPG-lite}: a lightweight DAPG-style demonstration-augmented policy-gradient
  baseline, abbreviated as \code{DAPG-lite} in the tables. It is included as a matched-protocol
  mechanism control rather than as a full
  reproduction of any single prior DAPG implementation.
\end{enumerate}

\begin{table}[t]
  \centering
  \small
  \caption{Operational definitions for the controlled factors used in the diagnostic sweeps.}
  \label{tab:factor_definitions}
  \begin{tabular}{ll}
    \toprule
    Term & Operational meaning in this paper \\
    \midrule
    Demonstration support & State-action support supplied by the BC demonstration dataset \\
    Train-reset coverage & Initial-state distribution used during BC/PPO training \\
    BC optimization depth & Number of supervised BC pretraining updates \\
    PPO budget & PPO environment interaction steps after optional BC warm-start \\
    \bottomrule
  \end{tabular}
\end{table}

\subsection{Support-Gated Variable-Impedance Learning and Support Coverage Index}

We refer to the main recipe studied in this paper as Support-Gated Variable-Impedance Learning
(SG-VI). In this repository SG-VI is not a general controller or a teacher-independent theorem. It
is the benchmark-local combination of three ingredients: (i) an explicit variable-impedance action
space, (ii) BC-based demonstration support that can reopen contact-relevant states before PPO
fine-tuning, and (iii) factorized controls that separate demonstration coverage, train-reset
coverage, BC optimization depth, and PPO budget.

To make the support claim measurable, we define a Support Coverage Index (SCI) over a projected
state-action signature
\begin{equation}
\psi(\mathbf{o}_t,\mathbf{a}_t) =
\left[
\|\mathbf{p}^{\mathrm{rel}}_{t,xy}\|_2,\;
p^{\mathrm{rel}}_{t,z},\;
\|\mathbf{f}_t\|_2,\;
\|[\Delta x_t,\Delta y_t]\|_2,\;
\Delta z_t,\;
\kappa^{xy}_t,\;
\kappa^{z}_t
\right].
\end{equation}
Let $Q_{\Delta}$ quantize this signature with benchmark-local bin widths
$\Delta = (0.5\,\mm,\; 0.5\,\mm,\; 0.25\,\N,\; 0.1,\; 0.1,\; 0.1,\; 0.1)$, and let
\begin{equation}
\mathcal{S}_D =
\left\{
Q_{\Delta}\left(\psi(\mathbf{o}^{D}_i,\mathbf{a}^{D}_i)\right)
\;\middle|\;
i = 1,\dots, |D|
\right\}
\end{equation}
be the quantized support set induced by a demonstration dataset $D$. For rollout $R$ with $T$
state-action pairs, SCI is
\begin{equation}
\mathrm{SCI}(R;D) = \frac{1}{T}\sum_{t=1}^{T}
\mathbb{I}\left[
Q_{\Delta}\left(\psi(\mathbf{o}^{R}_t,\mathbf{a}^{R}_t)\right) \in \mathcal{S}_D
\right].
\end{equation}
SCI is therefore 1 when the rollout remains entirely inside the quantized demonstration support and
0 when it never re-enters that support. In this paper SCI is proposed as a benchmark-local
diagnostic that complements success, contact, and force metrics rather than replacing them.

All learned baselines reported in the main text are trained on the nominal profile only. The stable imitation
configuration used throughout the main benchmark is
\code{bc\_rollout\_episodes = 32}, \code{bc\_pretrain\_steps = 32}, and
\code{bc\_batch\_size = 64}. This was selected because a separate stability sweep showed that
\code{32/32} is the smallest BC budget that remains stable across the critical seed combinations
(Appendix~\ref{app:bc-stability}).
We explicitly distinguish three notions that are easy to conflate: \emph{demonstration coverage}
is the state-action support present in the BC demonstrations, \emph{reset coverage} is the
initial-state support induced by the reset curriculum, and \emph{BC optimization stability} is
whether a given imitation budget reliably produces a successful BC base across seed combinations.
The \code{32/32} setting is used because it is the smallest budget that remains stable under the
current seed protocol, while Section~\ref{sec:coverage} manipulates reset coverage separately.
For the fixed-impedance learned baseline, stiffness adaptation is disabled by pinning the lateral
and axial stiffness ranges to constant values during both training and evaluation. Concretely, the
fixed baseline uses \mbox{$K_{xy}=65.0$} and \mbox{$K_z=87.5$}, i.e., the midpoint values of the
default lateral and axial stiffness ranges, and keeps those same constants in both training and
evaluation.
In addition to the learned comparison set, the evaluation package contains a minimal
same-interface handcrafted/classical anchor roster: \code{pose\_only}, a hand-crafted
\code{fixed\_impedance} controller, a hand-crafted \code{variable\_impedance} controller used to
generate demonstrations, and the added \code{hybrid\_position\_force},
\code{compliant\_search}, and \code{tuned\_impedance} references. We use them in the main text
only as benchmark-local engineering anchors under the shared interface; they improve external
reference coverage, but they are not presented as an exhaustive survey of classical insertion
controllers. Because the BC demonstrations are generated by the hand-crafted variable-impedance
teacher, the present design does not fully decouple demonstration support from the impedance prior
embedded in that teacher. The tuned fixed-stiffness sweep in Appendix~\ref{app:tuned-fixed}
partially checks whether fixed impedance can succeed when tuned, but it is not a matched-motion
or fixed-expert demonstration ablation.

\subsection{Nominal Profiles and Evaluation Classes}

The benchmark nominally evaluates five uncertainty profiles:
\code{nominal}, \code{tight\_clearance}, \code{high\_friction}, \code{offset\_bias}, and
\code{noisy\_force}. However, these do not correspond to five independent evaluation stresses in the
current environment semantics. Two design facts matter. First, \code{offset\_bias} is strongly
attenuated under the hole-centered relative observation frame used here. We do not yet include an
absolute-versus-relative observation ablation, so this should be read as an environment-semantic
explanation rather than as a separately isolated causal test. Second, \code{tight\_clearance} only
becomes active when the rollout enters lateral-overflow contact. For the reported discussion we
therefore group the five nominal profiles into three effective evaluation classes, referred to
below as evaluation classes:
\begin{itemize}
  \item \textbf{baseline}: \code{nominal}, \code{tight\_clearance}, and \code{offset\_bias},
  \item \textbf{high friction}: \code{high\_friction},
  \item \textbf{noisy force}: \code{noisy\_force}.
\end{itemize}

This is not a post-hoc averaging trick. It is a statement about environment semantics, and we make
it explicit because otherwise the repeated values across several profiles would look suspicious
rather than interpretable.

\subsection{Contact Model as a Precondition}

All benchmark results reported in the main text use the repaired analytical contact model as the default
environment. This is an intentional precondition rather than a hidden tuning choice. The original
and repaired contact-law variants are audited in Appendix~\ref{app:contact-audit}. In the main
text we keep the repaired contact model fixed so that demonstration support, train-reset
coverage, BC depth, and PPO budget can be compared under one protocol.

\subsection{Evaluation Contract}

The main benchmark uses five training seeds and 100 evaluation episodes per seed and profile.
The supplementary statistics report attaches bootstrap 95\% confidence intervals to the main
metrics and paired sign-permutation comparisons to the selected success-rate contrasts discussed
below. The main reported metrics are success rate, jam rate, mean final distance to target, mean
peak contact force, and mean contact steps. In the main tables, the evaluation-class success
columns report mean $\pm$ std across the five training seeds, and the aggregate columns report the
corresponding per-seed 5-profile means $\pm$ std. The factorized causal sweeps in
Section~\ref{sec:coverage} use a smaller diagnostic matched protocol and should therefore be read
as directional mechanism diagnostics rather than as the manuscript's highest-precision estimates.
Unless otherwise noted, all learned baselines reported in the main text also share the same
small PPO fine-tuning budget (\code{total\_timesteps = 128}) and the same episode horizon
(\code{max\_episode\_steps = 64}). Here \code{total\_timesteps} denotes PPO environment
interaction steps passed to the optimizer; it does not include BC rollout episodes or BC
pretraining updates. The large-budget PPO-only audit in Appendix~\ref{app:ppo-large-budget} is
reported separately because it changes the PPO-only training budget while preserving the same
evaluation profiles and episode horizon. We therefore report success gaps with their bootstrap
intervals and paired tests directly instead of assigning universal algorithm-ranking labels.
Throughout the results, effect size, interval, and paired test are treated as a joint description,
not as a single-p-value basis for ranking algorithms.

\section{Experiments}
\label{sec:experiments}

The experiments use separate evidence layers. Figure~\ref{fig:main_results_matrix} and
Table~\ref{tab:main_benchmark} report the final five-seed benchmark with 100 evaluation episodes
per seed and profile. The Sprint 2 reviewer-facing main table exported as
\code{outputs/evidence_matrix/three_dof_sprint2_main_table.md} serves a different role: it is a
three-layer claim-control table tying nominal-only pure-RL negatives, demo-supported
contact-reopening rows, and a mechanics / fixed-impedance anchor without turning them into a
single leaderboard. The cross-family PPO/SAC/TD3 pilot and its confirm report provide a separate
Branch-A negative-evidence layer, Table~\ref{tab:mechanism_matrix} and
Table~\ref{tab:behavioral_similarity} are smaller diagnostic sweeps intended for directional
mechanism analysis, the frozen Sprint 3 teacher mini-ablation kickoff fixes the next
teacher-support-quality $\times$ demo-budget question, and the Sprint 4A clearance sweep provides a
pure easy/nominal/hard clearance ladder that should not be read as a replacement for the frozen
five-profile benchmark. Figure~\ref{fig:high_friction_mechanism} reports high-friction mechanics
traces under the same 3DoF interface but with its own trace sampling contract,
Table~\ref{tab:pose_perturbation} reports a separate pose-perturbation proxy study, and Appendix
Tables~\ref{tab:ppo_large_budget} and~\ref{tab:tuned_fixed_sweep} report supplementary stress
checks rather than new main benchmark replacements.

\subsection{Main Benchmark: Coverage Matters More than Algorithm Choice}

Figure~\ref{fig:main_results_matrix} and Table~\ref{tab:main_benchmark} report the main
benchmark. Under the matched small-budget protocol, \code{PPO w/o BC} remains at 0.000 mean
success, and Appendix~\ref{app:ppo-large-budget} preserves the same non-contact failure at
50k--200k environment steps. Here ``useful contact'' means contact trajectories that enter the
recoverable contact regime and can still satisfy the benchmark success criteria, rather than merely
producing incidental force. Once the BC base is stabilized, \code{BC-only} reaches
$0.9996 \pm 0.0008$ mean success, while \code{BC$\rightarrow$PPO} and \code{DAPG-lite} both reach
$1.000 \pm 0.000$. The paired success gaps against \code{BC-only} are only 0.0004 with bootstrap CI
$[0.0000, 0.0012]$ and paired sign-permutation $p = 1.0$, so the final benchmark does not support a
meaningful success-rate ranking among the near-ceiling BC-based methods. Instead, it motivates the
factorized controls below, where demonstration support, reset coverage, BC depth, and PPO budget are
separated explicitly.

\begin{figure}[t]
  \centering
  \includeorplaceholder{fig2_main_benchmark_evaluation_class_summary}{Final Figure 2: main benchmark across evaluation classes}
  \caption{Main benchmark across evaluation classes. Each cell shows success rate for one
  method under one evaluation class, with the rightmost column reporting the 5-profile mean success.
  Under the matched small budget of this paper, PPO without BC warm-start does not reach useful contact
  behavior; Appendix~\ref{app:ppo-large-budget} shows that a PPO-only 50k--200k-step audit remains
  in the same non-contact regime. Meanwhile, \code{BC-only},
  \code{BC$\rightarrow$PPO}, and \code{DAPG-lite} remain inside a ceiling-saturated success band
  under the final five-seed statistical contract. Later factorized controls show that this similarity is
  budget-sensitive rather than universal. The
  fixed-impedance deficit is concentrated under \code{high\_friction}.}
  \label{fig:main_results_matrix}
\end{figure}

\begin{table}[t]
  \centering
  \small
  \caption{Main benchmark reported in the text. All learned methods are trained on the nominal profile and
  evaluated with 5 seeds and 100 episodes per seed per profile. The three evaluation-class columns
  report seed-level success mean $\pm$ std after grouping the five nominal profiles into
  \emph{baseline}, \emph{high friction}, and \emph{noisy force}. The remaining columns summarize
  the corresponding per-seed 5-profile mean behavior; bootstrap CI and paired-comparison notes are
  reported in the supplementary statistics report.}
  \label{tab:main_benchmark}
  \resizebox{\linewidth}{!}{%
  \begin{tabular}{lcccccccc}
    \toprule
    Method & Baseline & High fric. & Noisy force & Mean success & Jam & Final dist. (\mm) & Peak force (\N) & Contact steps \\
    \midrule
    PPO w/o BC & $0.00 \pm 0.00$ & $0.00 \pm 0.00$ & $0.00 \pm 0.00$ & $0.000 \pm 0.000$ & $0.000 \pm 0.000$ & $18.381 \pm 0.336$ & $0.000 \pm 0.000$ & $0.0 \pm 0.0$ \\
    BC-only & $1.00 \pm 0.00$ & $1.00 \pm 0.00$ & $1.00 \pm 0.00$ & $1.000 \pm 0.001$ & $0.000 \pm 0.000$ & $0.901 \pm 0.008$ & $0.932 \pm 0.033$ & $30.0 \pm 1.1$ \\
    Fixed-impedance RL & $0.99 \pm 0.02$ & $0.78 \pm 0.29$ & $0.98 \pm 0.04$ & $0.947 \pm 0.067$ & $0.000 \pm 0.000$ & $0.916 \pm 0.030$ & $1.071 \pm 0.109$ & $36.5 \pm 1.2$ \\
    BC$\rightarrow$PPO & $1.00 \pm 0.00$ & $1.00 \pm 0.00$ & $1.00 \pm 0.00$ & $1.000 \pm 0.000$ & $0.000 \pm 0.000$ & $0.902 \pm 0.007$ & $0.902 \pm 0.064$ & $29.8 \pm 1.5$ \\
    DAPG-lite & $1.00 \pm 0.00$ & $1.00 \pm 0.00$ & $1.00 \pm 0.00$ & $1.000 \pm 0.000$ & $0.000 \pm 0.000$ & $0.895 \pm 0.003$ & $1.009 \pm 0.101$ & $29.6 \pm 1.0$ \\
    \bottomrule
  \end{tabular}}
\end{table}

\subsection{Cross-Family Pure-RL Negative Evidence and Claim Control}

The manuscript's Branch-A narrative is locked by a nominal-train, nominal-eval cross-family pilot.
Across \code{PPO w/o BC}, \code{SAC w/o BC}, and \code{TD3 w/o BC} at 50k, 100k, and 200k
training budgets, none of the pure-RL runs enters useful contact: success remains 0, mean contact
steps remain 0, and the mean first-contact step stays at the 64-step episode horizon. The only
visible separation is terminal distance-to-contact, where \code{SAC w/o BC} is the strongest
distance proxy but still not a contact-reaching policy. To keep the cross-family negative evidence
tied to the main benchmark without collapsing them into
one leaderboard, we export a separate mixed-contract evidence matrix. Its role is reviewer-facing
claim control, not a new ranking table: the nominal-only pure-RL pilot rows
(\code{PPO w/o BC}, \code{SAC w/o BC}, \code{TD3 w/o BC}) remain explicit
\emph{contact-gate negatives}, while the five-profile benchmark rows
(\code{BC-only}, \code{BC$\rightarrow$PPO}, \code{DAPG-lite}, and the fixed-impedance learned anchor)
show that demonstration-supported methods reopen contact and non-zero success. In that matrix,
\code{SAC w/o BC} is retained only as the strongest \emph{distance proxy} among the pure-RL rows;
it is not a success baseline, and the mixed-contract matrix is not read as a direct leaderboard.
Each matrix row cites the direct confirm JSON or schema2 benchmark JSON passed to the exporter,
rather than a separate paper-table export, so the provenance stays inspectable at row level.
The same row contract is rendered as \code{outputs/evidence_matrix/three_dof_sprint2_main_table.md}.
Its companion pure-RL budget-curve figure,
\code{outputs/cross_family_confirm/three_dof_cross_family_confirm_learning_curve_summary.png},
plots distance proxy, success, and contact steps against the 50k--200k training budgets within the
nominal-only confirm contract. Thus \code{SAC w/o BC} appears only as a zero-contact distance
proxy and not a success baseline, solve-insertion claim, useful-contact claim, or cross-contract
winner.

The supplementary statistics report reinforces the same reading: within the ceiling-saturated
BC-based band, the informative differences are modest contact-load shifts rather than success-rate
separation.
\code{BC-only} and \code{BC$\rightarrow$PPO} remain close in mean peak contact force
($0.932 \pm 0.033$ versus $0.902 \pm 0.064$~\N), p95 peak contact force
($1.234 \pm 0.047$ versus $1.192 \pm 0.093$~\N), and mean contact steps
($30.0 \pm 1.1$ versus $29.8 \pm 1.5$). \code{DAPG-lite} carries slightly higher force statistics
($1.009 \pm 0.101$~\N mean peak force; $1.339 \pm 0.138$~\N p95) without any success advantage.
The fixed-impedance learned baseline sharpens the interpretation. Its overall mean success is
$0.947 \pm 0.067$, and the comparison against \code{BC$\rightarrow$PPO} has a larger success gap
(mean gap $-0.0532$, bootstrap CI $[-0.1208, -0.0064]$). At the same time, with only five paired
training seeds, the paired sign-permutation p-value is still $0.0625$, so this descriptive gap
should not be treated as a conventional $p<0.05$ null-hypothesis win. Its deficit is not uniform
across the evaluation classes: the main drop occurs under
\code{high\_friction}, where success falls to $0.776 \pm 0.287$. This localized failure mode sets
up the physical interpretation in Section~\ref{sec:physical-value}.

For additional same-interface engineering context, the simple hand-crafted references in the
repository show the same high-friction pattern. Under \code{high\_friction}, \code{pose\_only},
hand-crafted \code{fixed\_impedance}, and hand-crafted \code{variable\_impedance} all reach
1.000 success, but their mean peak contact forces differ sharply:
$1.58 \pm 0.00$~\N, $4.33 \pm 0.01$~\N, and $1.20 \pm 0.00$~\N, respectively. The repository also
keeps minimal same-interface classical anchors (\code{hybrid\_position\_force}, \code{compliant\_search}, and
\code{tuned\_impedance}) as supporting reference evidence. They strengthen the engineering
reference set inside the same benchmark interface, but they are not part of the final main
learned-method table and are not presented as an exhaustive classical-controller survey.

\subsection{Factorized Causal Study: Demonstration Support Reopens Contact, but Reset and Budget Effects Are Protocol-Sensitive}
\label{sec:coverage}

The main benchmark alone cannot distinguish among demonstration support, train-reset support,
imitation depth, and PPO budget. Table~\ref{tab:mechanism_matrix} therefore provides a
factorized diagnostic readout. Four findings matter.
First, removing demonstration support drives 5-profile mean success to 0.000, reduces mean
contact steps from 34.06 to 15.87, and raises mean jam from 0.105 to 0.452; the policy still
touches contact, but not in a useful way. Second, train-reset coverage also changes learnability,
but under the present diagnostic protocol its direction is not a universal monotone effect: the
narrow near-contact reset point reaches 1.000 mean success, whereas the broader reset point reaches
0.565. We therefore report reset coverage descriptively here rather than universalizing a broad
$>$ narrow claim. A plausible protocol-local explanation is that, at the small diagnostic budget,
the broader reset curriculum dilutes visits to contact-relevant near-hole states faster than the
BC/PPO budget can stabilize them; the narrow reset instead concentrates training on the recoverable
contact regime. Third, BC optimization depth has a clear sweet spot at 32 pretraining steps:
0 and 8 do not reopen success at all, 32 reaches 0.565 mean success, and 64 falls back to 0.456.
Finally, these controls show why the word ``coverage'' is too coarse on its own.
Demonstration support, train-reset support, and BC optimization stability do not move the task in
the same way, even under a shared contact model and matched evaluation contract.

\begin{table}[t]
  \centering
  \small
  \caption{Factorized causal controls under a smaller diagnostic matched protocol. Entries report
  5-profile mean success, mean contact steps, and mean jam rate for directional mechanism
  analysis; they are not direct numerical estimates under the final five-seed benchmark contract
  in Table~\ref{tab:main_benchmark}.}
  \label{tab:mechanism_matrix}
  \begin{tabular}{llccc}
    \toprule
    Factor & Condition & Mean success & Contact steps & Jam \\
    \midrule
    Demo support & removed & 0.000 & 15.87 & 0.452 \\
    Demo support & restored & 0.565 & 34.06 & 0.105 \\
    Train reset & narrow & 1.000 & 19.69 & 0.000 \\
    Train reset & broad & 0.565 & 34.06 & 0.105 \\
    BC pretrain steps & 0 & 0.000 & 0.00 & 0.000 \\
    BC pretrain steps & 8 & 0.000 & 11.72 & 0.317 \\
    BC pretrain steps & 32 & 0.565 & 34.06 & 0.105 \\
    BC pretrain steps & 64 & 0.456 & 32.78 & 0.027 \\
    \bottomrule
  \end{tabular}
\end{table}

This result is important for the paper's causal story. A single reset-coverage narrative would
compress too many distinct mechanisms. The factorized controls show a cleaner but more qualified
picture: demonstration support is the robust gating variable, train-reset coverage is real but
protocol-local, and BC optimization depth matters as a non-monotone stability knob.

\subsection{High-Friction Role of Variable Impedance}
\label{sec:physical-value}

The fixed-versus-variable comparison allows a separate question: if algorithm choice is largely
secondary once BC is stable, where does variable impedance still matter physically? The answer is
not in the baseline evaluation class, where fixed impedance remains strong, nor mainly in
\code{noisy\_force}, where the gap is modest. The main separation occurs under
\code{high\_friction}. Under the final five-seed benchmark, the stable fixed-impedance
learned baseline reaches $0.776 \pm 0.287$ success under the same nominal training profile,
whereas the variable-impedance pipelines remain at 1.000 success.
The same-interface engineering anchors and the tuned fixed-stiffness sweep in
Appendix~\ref{app:tuned-fixed} indicate that fixed impedance is not fundamentally incapable of
success in this reduced task. The sweep selects $K_{xy}=80$ and $K_z=50$ and confirms a
5-profile mean success of 0.999, including 1.000 under \code{high\_friction}. This removes the
stronger interpretation that fixed impedance only fails because the midpoint stiffness is untuned.
The remaining fixed-versus-variable claim is therefore narrower: compared with
variable impedance, fixed impedance can succeed when tuned, but the main learned midpoint-fixed
pipeline is less robust under the matched training budget, and the same-interface mechanics traces
show that fixed strategies can operate at substantially higher contact load.

Figure~\ref{fig:high_friction_mechanism} is a diagnostic mechanics slice, not a replot of the
main five-seed benchmark. It aggregates high-friction trace rollouts from three training seeds and
50 trace episodes per seed for learned fixed and learned variable suites, together with
same-interface hand-crafted fixed and variable anchors. The hand-crafted fixed anchor uses the
policy's default fixed command, which decodes to approximately $K_{xy}=74$ and $K_z=98$ under the
default stiffness range; it is not the success-optimized fixed setting selected in
Appendix~\ref{app:tuned-fixed}. The learned traces are intentionally shown
as heterogeneous all-trace aggregates; their success rates should therefore not be compared
numerically with Table~\ref{tab:main_benchmark}. The cleanest mechanics contrast comes from the
two hand-crafted anchors because both succeed reliably under this diagnostic trace protocol while
operating at different contact loads. By the end of the aggregated contact phase, the hand-crafted
fixed controller remains near $0.87$~\N\ total contact force and about $1.87\times 10^{-2}$~J
cumulative contact work, whereas the hand-crafted variable controller remains near $0.46$~\N\ and
about $4.35\times 10^{-3}$~J. This direct evidence supports a conservative interpretation: fixed
impedance is not physically impossible in this reduced task, but variable impedance can realize a
lower-load, lower-work contact path under high friction. The learned all-trace view remains
heterogeneous, which is why the released figure family also includes a successful-only secondary
aggregation rather than claiming a single clean monotone force ranking from the learned traces
alone.

\begin{figure}[t]
  \centering
  \includeorplaceholder{fig3_high_friction_impedance_mechanism}{Final Figure 3: high-friction impedance mechanism}
  \caption{Diagnostic high-friction mechanics traces under the shared 3DoF benchmark interface.
  This figure is a trace-level mechanics slice, not the final five-seed benchmark table: the
  learned fixed and learned variable traces use three training seeds with 50 trace episodes per
  seed, while the hand-crafted anchors use the same trace count. The all-trace learned aggregates
  are heterogeneous and are included to show rollout mechanics rather than to provide success-rate
  estimates comparable to Table~\ref{tab:main_benchmark}. The cleanest physical contrast is the
  hand-crafted fixed-vs-variable pair, where both controllers succeed reliably but the variable
  strategy follows a lower-load, lower-work contact path. The hand-crafted fixed anchor is not the
  tuned fixed-stiffness setting from Appendix~\ref{app:tuned-fixed}; it is a same-interface
  mechanics reference for fixed-command contact paths. A successful-only companion view is provided
  in the supplementary material.}
  \label{fig:high_friction_mechanism}
\end{figure}

\subsection{Supplementary Budget Stress Test After BC Support Is Restored}

The diagnostic budget comparison further narrows the algorithm claim. \code{PPO w/o BC} remains
at 0.000 mean success at PPO budgets of 0, 32, 128, and 512 timesteps, so increasing PPO budget
within this small sweep does not recover the task. The larger PPO-only audit in
Appendix~\ref{app:ppo-large-budget} extends this negative result to 50k, 100k, and 200k PPO
environment steps and shows zero contact steps and zero peak contact force across both PPO-only
training contracts. \code{BC$\rightarrow$PPO} is not monotone
either: its 5-profile
mean success drops from 0.675 at zero extra PPO steps to 0.644 at 32, 0.565 at 128, and 0.012 at
512. \code{DAPG-lite} is also non-monotone, reaching 0.673, 0.061, 0.620, and 0.325 at the same
budgets. The static \code{BC-only} anchor remains at 1.000. Table~\ref{tab:behavioral_similarity}
therefore no longer supports a ``once BC is stable, all downstream algorithms saturate'' reading.
The safer interpretation is that under the matched protocol of this paper, algorithm ranking is
local to the available demonstration support, BC optimization depth, and PPO budget. BC-based
support is still the gating condition, but PPO fine-tuning is not a guaranteed monotone gain and
can regress sharply at larger budgets.

\begin{table}[t]
  \centering
  \small
  \caption{Diagnostic algorithm comparison as PPO budget varies under restored demonstration
  support and broad train-reset coverage. Entries are 5-profile mean success in the smaller sweep
  and are intended to show budget sensitivity rather than to replace the final five-seed benchmark
  in Table~\ref{tab:main_benchmark}.}
  \label{tab:behavioral_similarity}
  \begin{tabular}{lcccc}
    \toprule
    Method & 0 & 32 & 128 & 512 \\
    \midrule
    PPO w/o BC & 0.000 & 0.000 & 0.000 & 0.000 \\
    BC$\rightarrow$PPO & 0.675 & 0.644 & 0.565 & 0.012 \\
    DAPG-lite & 0.673 & 0.061 & 0.620 & 0.325 \\
    BC-only anchor & \multicolumn{4}{c}{1.000 (static)} \\
    \bottomrule
  \end{tabular}
\end{table}

\subsection{Teacher Support Quality vs.\ Demo Budget Boundary}

The next teacher-coupling question is kept intentionally narrow. Rather than claiming a finished
teacher ablation from partial runs, we freeze a 2 $\times$ 2 boundary over teacher support quality
(\emph{support rich} versus \emph{support poor}) and demonstration rollout budget
(\emph{many demo} versus \emph{few demo}). This boundary keeps BC pretraining steps, PPO
fine-tuning steps, policy initialization, and evaluation metrics fixed while varying only the two
teacher-coupled axes that matter for the next iteration. In other words, the current Sprint 3
artifact is a pre-registered boundary condition, not a new leaderboard or an additional benchmark
result. Its role in the paper is methodological: it prevents us from expanding the teacher story
post hoc after seeing partial training outcomes, and it keeps the eventual teacher ablation tied to
contact-gate contrast plus support diagnostics rather than to a broad teacher-independent claim.

\subsection{Clearance-Shift Robustness Sweep}

The frozen five nominal profiles do not isolate pure clearance. Under the current environment
semantics, \code{tight\_clearance} becomes active only under lateral-overflow contact, and
\code{offset\_bias} is strongly attenuated by the hole-centered relative frame. Sprint 4A therefore
adds a separate pure-clearance ladder with three evaluation profiles:
\code{clearance\_easy} with a 0.95--1.35~\mm\ clearance range, the nominal 0.70--1.10~\mm\
benchmark range, and \code{clearance\_hard} with a 0.45--0.75~\mm\ range. This ladder keeps hole
offset, friction, force noise, and the train profile fixed, so it is a narrower stress test than
the five-profile benchmark rather than a replacement for it.

Under this ladder, \code{BC-only} remains at 1.000 success across all three clearance levels.
\code{BC$\rightarrow$PPO} stays near 0.8 success and shows no jam increase, while the learned
fixed-impedance baseline also remains near 0.8 success without jam. \code{DAPG-lite} is the most
sensitive selected suite: it drops from 0.800 success under \code{clearance\_easy} to 0.768 under
\code{clearance\_hard}, while jam rises from 0.000 to 0.050 and the mean final distance grows from
1.178 to 1.617~\mm. We therefore read the clearance ladder conservatively. It does not overturn
the frozen five-profile benchmark, but it does show that harder pure-clearance shifts are not
equally benign for all demonstration-supported pipelines.

\subsection{Pose-Perturbation Scope Stress Test}

The benchmark is still analytical and simulation-only, so any external-validity extension must be
read conservatively. We therefore add one explicit proxy study rather than pretending that the
current evidence has full 6D generalization. The dedicated pose-perturbation study defines three
same-interface evaluation profiles with hidden contact-target bias and descent-coupled lateral
drift: \code{pose\_perturb\_mild} uses a 0.2~\mm\ hidden lateral bias with a 1.5$^\circ$ descent
coupling proxy, \code{pose\_perturb\_moderate} uses 0.4~\mm\ and 3.0$^\circ$, and
\code{pose\_perturb\_severe} uses 0.6~\mm\ and 4.5$^\circ$. The key point is not that this proxy
recreates full orientation dynamics. It does not. Its value is that it breaks the clean
hole-centered relative-frame assumption in a controlled, auditable way while keeping the same 3DoF
interface and the same benchmark metrics.

The result is a scope stress test rather than a second main benchmark. The mild proxy leaves the
stable \code{BC-only} suite near ceiling (1.000 success), but the same hidden perturbation already drops
\code{BC$\rightarrow$PPO} to 0.480 and \code{DAPG-lite} to 0.446. Under the moderate proxy,
\code{BC-only} remains at 0.946 while the PPO-tuned learned suites stay near 0.39--0.45 and the
stable fixed-impedance learned baseline drops to 0.784. Under the severe proxy, all three learned
variable-impedance pipelines degrade sharply: \code{BC-only} reaches 0.302, \code{BC$\rightarrow$PPO}
reaches 0.306, and \code{DAPG-lite} reaches 0.320, while the stable fixed-impedance learned
baseline drops further to 0.110. In contrast, the same-interface
\code{tuned\_impedance} handcrafted anchor remains at 1.000 across all three proxy levels.
This means the main matched-protocol ranking does not survive unchanged once the relative-frame
assumption is stressed. The benchmark still teaches something useful, but its scope stays
benchmark-local even after the dedicated pose-perturbation study.

\begin{table}[t]
  \centering
  \small
  \caption{Separate pose-perturbation proxy study under the shared 3DoF interface. Entries report
  success rate under three hidden pose/tilt proxy levels. This is a scope-stressing proxy, not a
  full 6D orientation or hardware validation, and it is not directly comparable to the matched
  main benchmark.}
  \label{tab:pose_perturbation}
  \begin{tabular}{lccc}
    \toprule
    Method & Mild & Moderate & Severe \\
    \midrule
    BC-only & 1.000 & 0.946 & 0.302 \\
    Fixed-impedance RL & 0.988 & 0.784 & 0.110 \\
    BC$\rightarrow$PPO & 0.480 & 0.390 & 0.306 \\
    DAPG-lite & 0.446 & 0.454 & 0.320 \\
    Hand-tuned impedance & 1.000 & 1.000 & 1.000 \\
    \bottomrule
  \end{tabular}
\end{table}

\section{Discussion and Limitations}
\label{sec:discussion}

Framed as SG-VI, the main lesson of the benchmark is not that a particular pure-RL optimizer can
never solve insertion in principle. It is that under the present interface, pure RL across PPO,
SAC, and TD3 does not reach contact-relevant behavior before upstream support becomes the dominant
issue. The dedicated PPO-only audit extends that negative result to 200k environment steps under
two PPO contracts, while the cross-family pilot shows that changing optimizer family alone still
does not reopen useful contact. Demonstrations therefore act here as a learnability scaffold rather
than as a mere sample-efficiency convenience.

The second lesson is that the upstream interface must be split carefully. Demonstration coverage
is the clean one-directional lever: removing it drives useful-contact success to zero. Train-reset
coverage also changes learnability, but in the current matched protocol its direction is
protocol-local rather than universal or monotone. BC optimization depth and PPO budget are likewise
non-monotone; 32 BC pretraining steps is the current sweet spot, and additional PPO steps can
substantially hurt a BC-initialized policy. This is why we instead separate demonstration coverage,
reset coverage, BC optimization stability, and PPO budget instead of treating ``coverage'' or
``algorithm choice'' as single knobs. SCI is meant to accompany that decomposition rather than
collapse it into a single leaderboard scalar: it asks whether a rollout stays inside the available
demonstration support before asking how successful that rollout eventually becomes. The frozen
teacher-support-quality $\times$ demo-budget boundary serves the same purpose methodologically: it
keeps the next teacher ablation auditable instead of letting the teacher story expand post hoc.

The third lesson is that variable impedance still matters, but in a physically localized way.
Fixed impedance is not uniformly poor in this benchmark: the tuned fixed-stiffness sweep nearly
saturates success, including under high friction. The main fixed-impedance learned gap is therefore
not a broad impossibility result. It is a narrower robustness and load-path result: the midpoint
fixed learned pipeline is most fragile under high-friction contact, and the mechanics traces show
that variable stiffness can realize lower-load contact paths.

The fourth lesson is that robustness outside the frozen five-profile benchmark is axis-dependent.
The pure-clearance ladder does not collapse every selected demo-supported suite: \code{BC-only}
stays at 1.000 success, \code{BC$\rightarrow$PPO} and the learned fixed-impedance baseline remain
near 0.8 success without jam, and the clearest harder-clearance penalty appears in
\code{DAPG-lite}. By contrast, the dedicated pose-perturbation study changes the learned ranking
much more aggressively once hidden pose bias and descent-coupled lateral drift are introduced. That
is why the paper still keeps \code{benchmark-local} language even after adding both the
clearance-only and external-validity proxy stress tests.

Several limitations remain important. The environment is still simulation-only, translational, and
3DoF. It uses a relative observation frame, a synthetic force signal, diagonal stiffness decoding,
and an analytical contact model rather than torque sensing, 6D pose error, multi-point collision,
or hardware force/torque calibration. The pose-perturbation study improves stress coverage, but it
is still a proxy study rather than a hardware validation or a full 6D orientation benchmark. The
current statements about force-control accuracy, force-pose coupling, and compliant load paths
must therefore remain benchmark-local. The present project is also explicitly on the no-hardware
path, so the paper currently has no real-robot or cross-simulator transfer evidence to lift that
scope ceiling.

The learning evidence also leaves several deliberate separations unresolved. The large-budget
PPO-only audit covers two PPO contracts rather than an exhaustive hyperparameter search, so it
supports the benchmark-local non-contact finding without proving that every possible PPO variant
must fail. The BC demonstrations are generated by a variable-impedance teacher, which means
demonstration support and variable-impedance prior are not fully decoupled. A fixed-impedance
expert demonstration set, a matched-motion/mismatched-impedance study, or a fixed-expert BC
pipeline would be needed to isolate that prior more completely. The same-interface
handcrafted/classical anchors improve reference coverage, but the repository still does not
provide an exhaustive survey of hybrid position/force, admittance/compliance, demo-augmented RL,
or offline RL baselines. Accordingly, both SG-VI and SCI remain benchmark-local, teacher-coupled
definitions in the present paper rather than universal insertion metrics.

Finally, the metric and parameter sensitivity evidence is still narrower than a full systems
validation. The main table reports success, jam, final distance, peak force, and contact steps;
contact work and stiffness-progress traces remain diagnostic rather than primary Pareto metrics.
Future versions should add force-tracking error, contact work, success-force Pareto summaries, and
distance-energy joint metrics as first-class table entries. Likewise, \code{offset\_bias} is only
interpreted through the hole-centered relative-frame semantics; we do not yet provide an explicit
absolute-versus-relative observation ablation. The analytical contact parameters
$(\alpha_{xy}, \alpha_z, \mu_{\mathrm{wall}})$, jam thresholds, stiffness ranges, and demonstration
quantity/quality also need systematic sensitivity sweeps before the conclusions can be read as
anything beyond controlled-benchmark evidence. The same applies to SCI itself: the current
projected state-action signature and bin widths are frozen to the scale of this 3DoF relative-frame
benchmark rather than tuned as a universal discretization. The $0.5\,\mm$ position bins resolve
sub-millimeter contact-entry structure relative to the $0.8/1.0\,\mm$ success tolerances, the
$0.25\,\N$ force bin stays well below the $8\,\N$ jam threshold while separating low-force approach
from contact, and the $0.1$ normalized action bins keep the explicit motion/stiffness channels
coarse enough to avoid near-exact sample matching. These choices remain benchmark-local and still
require sensitivity analysis before SCI should be read as anything broader. A larger
10-seed expansion would tighten the uncertainty on the remaining near-ceiling gaps, but it would
not by itself close these external validity questions. Likewise, the Sprint 4A clearance sweep is
currently executed as a train-once / eval-many-profile stress test because this repository does not
yet persist paper-ready checkpoints for true eval-only reruns.

\section{Conclusion}
\label{sec:conclusion}

This paper presents Support-Gated Variable-Impedance Learning (SG-VI) as a benchmark-local way to
study 3DoF variable-impedance insertion when the main question is not which RL algorithm is
strongest, but which upstream conditions make the task learnable at all. Support Coverage Index
(SCI) provides a concrete rollout-to-demo support diagnostic for that question. Under the matched
protocol used here, PPO without BC warm-start does not reach useful contact, and a 50k--200k-step
PPO-only audit preserves the same non-contact failure. More precisely, BC-based demonstration
support is necessary to reach the useful contact regime, train-reset coverage remains a
protocol-local rather than universal lever, and downstream algorithm ranking is budget-sensitive
rather than fixed once PPO fine-tuning is varied. The learned fixed-impedance gap is concentrated
under high-friction contact, while the tuned fixed-stiffness sweep shows that fixed impedance can
nearly saturate success when tuned. This supports a localized load-efficiency role and a more
robust learned interface under the matched teacher-coupled setup rather than a diffuse claim of
uniform superiority. The dedicated pose-perturbation study further shows that this matched-frame
ranking does not automatically survive outside the clean hole-centered frame: several learned suites
degrade sharply under hidden pose bias and descent-coupled lateral drift, while the strongest
same-interface handcrafted anchor remains robust. Because the BC demonstrations are generated by a
variable-impedance teacher, these results identify a learnability gate under this teacher-coupled
interface rather than a teacher-independent causal theorem about all demonstration designs. These
conclusions are therefore benchmark-local and should be interpreted together with the simulation-only
3DoF scope, the proxy nature of the perturbation extension, and the still-limited classical anchor
roster.

\FloatBarrier
\appendix

\renewcommand{\thefigure}{A\arabic{figure}}
\setcounter{figure}{0}
\renewcommand{\thetable}{A\arabic{table}}
\setcounter{table}{0}

\section{Protocol Summary}
\label{app:protocol-summary}

Table~\ref{tab:protocol_summary} summarizes which evidence blocks share a numerical contract and
which should be read only directionally. This table is included to prevent direct comparison of
diagnostic sweeps or trace slices with the final five-seed benchmark.

\begin{table}[!htbp]
  \centering
  \small
  \caption{Protocol map for the paper's evidence blocks.}
  \label{tab:protocol_summary}
  \resizebox{\linewidth}{!}{%
  \begin{tabular}{lllll}
    \toprule
    Evidence block & Sampling contract & Train profile & Evaluation profiles & Role \\
    \midrule
    Table~\ref{tab:main_benchmark} / Figure~\ref{fig:main_results_matrix} & 5 seeds, 100 episodes per seed/profile & nominal & five nominal profiles grouped into three evaluation classes & final benchmark estimate \\
    Cross-family pure-RL pilot & 3 seeds, 50 episodes per seed/profile & nominal & nominal only & Branch-A negative evidence for PPO/SAC/TD3, not a mixed-contract leaderboard \\
    Teacher mini-ablation kickoff boundary & 4 frozen conditions, 5-profile evaluation & nominal & five nominal profiles & pre-registered teacher-support-quality $\times$ demo-budget boundary, not a completed ablation \\
    Sprint 4A clearance shift sweep & 5 seeds, 100 episodes per seed/profile & nominal & \code{clearance\_easy} / \code{nominal} / \code{clearance\_hard} & pure-clearance robustness stress, not a replacement for the frozen five-profile benchmark \\
    Table~\ref{tab:mechanism_matrix} & smaller diagnostic sweep & nominal & five nominal profiles & directional factor analysis, not direct comparison with Table~\ref{tab:main_benchmark} \\
    Table~\ref{tab:behavioral_similarity} & smaller diagnostic sweep & nominal & five nominal profiles & PPO-budget sensitivity, not direct comparison with Table~\ref{tab:main_benchmark} \\
    Figure~\ref{fig:high_friction_mechanism} & 3 seeds, 50 trace episodes per seed & nominal & high friction & mechanics trace slice, not a success estimate comparable to Table~\ref{tab:main_benchmark} \\
    Table~\ref{tab:pose_perturbation} & separate proxy study & nominal & hidden pose/tilt proxy levels & scope stress test, not a main benchmark replacement \\
    Table~\ref{tab:ppo_large_budget} & 5 seeds, 100 episodes per seed/profile & nominal & five nominal profiles & PPO-only budget audit, not a BC-method comparison \\
    Table~\ref{tab:tuned_fixed_sweep} & 3-seed search, 5-seed confirm & nominal & five nominal profiles & tuned fixed-stiffness check, not a variable-impedance mechanics substitute \\
    \bottomrule
  \end{tabular}}
\end{table}
\FloatBarrier

\section{Large-Budget PPO-Only Audit}
\label{app:ppo-large-budget}

The main benchmark uses a deliberately small matched PPO budget so that all learned methods share
one evaluation contract. To address the separate concern that PPO from scratch may simply need a
larger budget, we ran a PPO-only audit with BC rollout episodes and BC pretraining steps both set
to zero. The audit uses the same five evaluation profiles, 5 training seeds, 100 evaluation
episodes per seed and profile, and a 64-step episode horizon. It includes two PPO-only contracts:
the paper-matched optimizer settings and an alternative PPO setting with 4 parallel
environments, 256 rollout steps, 4 epochs, learning rate $3\times 10^{-4}$, $\gamma=0.99$, and
entropy coefficient 0.01. Table~\ref{tab:ppo_large_budget} shows that neither contract reaches
contact-relevant behavior up to 200k environment steps. A first-contact step equal to 64 indicates
that first contact was not observed before the episode horizon.

\begin{table}[!htbp]
  \centering
  \small
  \caption{PPO-only large-budget audit. Entries are 5-profile means over 5 training seeds and
  100 evaluation episodes per seed/profile. Both PPO-only contracts keep BC rollout episodes and
  BC pretraining steps at zero.}
  \label{tab:ppo_large_budget}
  \resizebox{\linewidth}{!}{%
  \begin{tabular}{llccccc}
    \toprule
    PPO-only contract & PPO steps & Success & Contact steps & First contact step & Final dist. (\mm) & Peak force (\N) \\
    \midrule
    Paper-matched & 50k & 0.000 & 0.0 & 64.0 & 21.59 & 0.000 \\
    Paper-matched & 100k & 0.000 & 0.0 & 64.0 & 23.96 & 0.000 \\
    Paper-matched & 200k & 0.000 & 0.0 & 64.0 & 24.42 & 0.000 \\
    Alternative PPO & 50k & 0.000 & 0.0 & 64.0 & 24.70 & 0.000 \\
    Alternative PPO & 100k & 0.000 & 0.0 & 64.0 & 25.04 & 0.000 \\
    Alternative PPO & 200k & 0.000 & 0.0 & 64.0 & 23.42 & 0.000 \\
    \bottomrule
  \end{tabular}}
\end{table}
\FloatBarrier

\section{Tuned Fixed-Impedance Sweep}
\label{app:tuned-fixed}

The main fixed-impedance learned baseline pins the impedance decoder to the midpoint values
$K_{xy}=65$ and $K_z=87.5$. To check whether this midpoint choice creates an untuned fixed
baseline, we ran a same-interface grid search over
$K_{xy}\in\{30,50,65,80,100\}$ and $K_z\in\{50,75,87.5,110,130\}$ using the learned fixed
pipeline with 32 BC rollout episodes, 32 BC pretraining steps, 128 PPO steps, 3 search seeds, and
50 evaluation episodes per seed/profile. Candidates were ranked by 5-profile mean success,
high-friction success, mean peak contact force, and mean final distance. The selected fixed
stiffness was then confirmed with 5 seeds and 100 evaluation episodes per seed/profile.
Table~\ref{tab:tuned_fixed_sweep} shows that tuned fixed impedance can nearly saturate success in
this reduced task. This check limits the fixed-versus-variable claim: fixed impedance is not a
success-rate impossibility result, so the variable-impedance interpretation should rest on the more
robust learned interface under the matched teacher-coupled setup and on the lower-load mechanics
evidence in Figure~\ref{fig:high_friction_mechanism}.

\begin{table}[!htbp]
  \centering
  \small
  \caption{Tuned fixed-impedance confirmation for the selected grid point. Entries are mean success
  over 5 seeds and 100 evaluation episodes per seed/profile.}
  \label{tab:tuned_fixed_sweep}
  \begin{tabular}{lccc}
    \toprule
    Profile & Success & Peak force (\N) & Final dist. (\mm) \\
    \midrule
    Nominal & 1.000 & 0.582 & 0.889 \\
    Tight clearance & 1.000 & 0.567 & 0.890 \\
    High friction & 1.000 & 0.920 & 0.904 \\
    Offset bias & 1.000 & 0.568 & 0.895 \\
    Noisy force & 0.996 & 0.776 & 0.906 \\
    5-profile mean & 0.999 & 0.683 & 0.897 \\
    \bottomrule
  \end{tabular}
\end{table}
\FloatBarrier

\section{Evaluation-Class Transparency}
\label{app:evaluation-classes}

Figure~\ref{fig:evaluation_class_mapping} summarizes why five nominal uncertainty profiles collapse
into only three evaluation classes in the current 3DoF environment. The purpose of this
appendix figure is not to add a fourth result, but to make the environment semantics explicit and
auditable. Table~\ref{tab:profile_mapping} states the same mapping in tabular form so that the
grouping logic is inspectable without reading the figure.

\begin{figure}[!htbp]
  \centering
  \includeorplaceholder{figA1_evaluation_class_mapping}{Appendix Figure A1: evaluation-class mapping}
  \caption{Mapping from the five nominal uncertainty profiles to the three evaluation
  classes used in the reported result. \code{nominal}, \code{tight\_clearance}, and
  \code{offset\_bias} form the baseline class because tight clearance is active only under lateral
  overflow and offset bias is strongly attenuated under the relative observation frame.
  \code{high\_friction} and
  \code{noisy\_force} remain separate because they alter contact mechanics and sensing uncertainty,
  respectively.}
  \label{fig:evaluation_class_mapping}
\end{figure}
\FloatBarrier

\begin{table}[!htbp]
  \centering
  \small
  \caption{Mapping from nominal profiles to evaluation classes used in the main paper.}
  \label{tab:profile_mapping}
  \begin{tabular}{lll}
    \toprule
    Nominal profile & Effective class & Rationale \\
    \midrule
    Nominal & Baseline & No extra active evaluation stress under current semantics \\
    Tight clearance & Baseline & Active only under lateral overflow \\
    Offset bias & Baseline & Strongly attenuated under the relative observation frame \\
    High friction & High friction & Changes contact mechanics \\
    Noisy force & Noisy force & Changes sensing uncertainty \\
    \bottomrule
  \end{tabular}
\end{table}
\FloatBarrier

\section{Per-Profile Breakdown}
\label{app:per-profile}

Table~\ref{tab:raw_profile_breakdown} complements the grouped main results by reporting the raw
five-profile values directly. Each entry is shown as \emph{success / jam}, so the appendix can be
used to verify both the grouped evaluation-class means in Table~\ref{tab:main_benchmark} and the
statement that jam remains absent throughout the main five-suite comparison.

\begin{table}[!htbp]
  \centering
  \small
  \caption{Raw nominal-profile breakdown for the main five learned suites. Entries are reported as
  \emph{success / jam}; the rightmost column reports the 5-profile mean in the same format.}
  \label{tab:raw_profile_breakdown}
  \resizebox{\linewidth}{!}{%
  \begin{tabular}{lcccccc}
    \toprule
    Method & Nominal & Tight clr. & High fric. & Offset bias & Noisy force & 5-profile mean \\
    \midrule
    PPO w/o BC & 0.000 / 0.000 & 0.000 / 0.000 & 0.000 / 0.000 & 0.000 / 0.000 & 0.000 / 0.000 & 0.000 / 0.000 \\
    BC-only & 1.000 / 0.000 & 1.000 / 0.000 & 1.000 / 0.000 & 1.000 / 0.000 & 0.998 / 0.000 & 1.000 / 0.000 \\
    Fixed-impedance RL & 0.994 / 0.000 & 0.988 / 0.000 & 0.776 / 0.000 & 0.992 / 0.000 & 0.984 / 0.000 & 0.947 / 0.000 \\
    BC$\rightarrow$PPO & 1.000 / 0.000 & 1.000 / 0.000 & 1.000 / 0.000 & 1.000 / 0.000 & 1.000 / 0.000 & 1.000 / 0.000 \\
    DAPG-lite & 1.000 / 0.000 & 1.000 / 0.000 & 1.000 / 0.000 & 1.000 / 0.000 & 1.000 / 0.000 & 1.000 / 0.000 \\
    \bottomrule
  \end{tabular}}
\end{table}
\FloatBarrier

\section{Contact-Transition Audit}
\label{app:contact-audit}

The repaired contact model used in the reported benchmark was not selected arbitrarily. The
audit in Figure~\ref{fig:contact_audit} shows that the repaired transition broadens the usable
safe-contact window and delays the onset of first jam. This is why the repaired model is treated
as a precondition for the main benchmark. The old-versus-new contact-law comparison is retained
only as an appendix audit that justifies the repaired contact model used throughout the main
benchmark.

\begin{figure}[!htbp]
  \centering
  \includeorplaceholder{fig1_contact_transition_audit}{Appendix Figure A2: contact-transition audit}
  \caption{Original versus repaired contact-transition audit. The repaired model broadens the
  safe-contact window from roughly 0.85 to 1.25~\mm\ and shifts the first-jam point from roughly
  0.90 to 1.30~\mm. This appendix figure justifies the repaired contact model as the default
  reported precondition without reinstating a dual-bottleneck claim.}
  \label{fig:contact_audit}
\end{figure}
\FloatBarrier

\section{BC Stability and Budget Evidence}
\label{app:bc-stability}

The main paper uses \code{bc\_rollout\_episodes = 32}, \code{bc\_pretrain\_steps = 32}, and
\code{bc\_batch\_size = 64} because this is the smallest BC budget that remains stable across the
critical seed combinations examined in the sweep. Table~\ref{tab:bc_stability} summarizes the key
comparison. The unstable point is not a vague degradation: increasing the rollout count to
\code{64} while keeping pretraining at \code{32} collapses the mean success to 0.352 and wipes out
the \code{noisy\_force} class entirely. Increasing optimization depth to \code{128} recovers the
performance, which is exactly the data pattern behind the paper's ``stable BC base'' wording.

\begin{table}[!htbp]
  \centering
  \small
  \caption{BC stability sweep for the critical seed combinations. All settings use
  \code{bc\_batch\_size = 64}.}
  \label{tab:bc_stability}
  \begin{tabular}{lcccccc}
    \toprule
    Setting & Init/demo & Mean success & High fric. & Noisy force & Final dist. (\mm) & Peak force (\N) \\
    \midrule
    r32 / p32 & 1 / 0 & 1.000 & 1.000 & 1.000 & 0.895 & 1.063 \\
    r32 / p32 & 1 / 1 & 1.000 & 1.000 & 1.000 & 0.896 & 0.950 \\
    r64 / p32 & 1 / 1 & 0.352 & 0.040 & 0.000 & 1.388 & 0.878 \\
    r64 / p128 & 1 / 1 & 1.000 & 1.000 & 1.000 & 0.920 & 0.969 \\
    \bottomrule
  \end{tabular}
\end{table}
\FloatBarrier

\bibliographystyle{plainnat}
\bibliography{references}

\end{document}
